\documentclass[10pt,twocolumn]{article}

\usepackage[margin=1in]{geometry}
\usepackage{amsmath,amssymb,amsfonts,amsthm}
\usepackage{mathtools}
\usepackage{graphicx}
\usepackage{hyperref}
\usepackage{xcolor}
\usepackage{booktabs}
\usepackage{array}
\usepackage{algorithm}
\usepackage{algpseudocode}
\usepackage[numbers,sort&compress]{natbib}
\usepackage{enumitem}

% Theorem environments
\newtheorem{theorem}{Theorem}[section]
\newtheorem{lemma}[theorem]{Lemma}
\newtheorem{corollary}[theorem]{Corollary}
\newtheorem{definition}[theorem]{Definition}
\newtheorem{proposition}[theorem]{Proposition}
\newtheorem{axiom}{Axiom}
\newtheorem{principle}[theorem]{Principle}

\theoremstyle{remark}
\newtheorem{remark}[theorem]{Remark}

% Custom commands
\newcommand{\Sk}{S_k}
\newcommand{\St}{S_t}
\newcommand{\Se}{S_e}
\newcommand{\Sspace}{\mathcal{S}}
\newcommand{\Scoord}{\mathbf{S}}
\newcommand{\Pcat}{\mathcal{P}}
\newcommand{\Tcat}{\mathcal{T}}
\newcommand{\Ccat}{\mathcal{C}}
\newcommand{\Pop}{P}
\newcommand{\Gstate}{\Gamma}
\newcommand{\kB}{k_{\mathrm{B}}}
\newcommand{\Hfield}{\mathcal{H}^+}
\newcommand{\Pdecay}{P_{\mathrm{decay}}}
\newcommand{\Tdecay}{T_{\mathrm{decay}}}
\newcommand{\Mstate}{\mathcal{M}}
\newcommand{\dS}{d_S}
\newcommand{\Sosc}{S_{\mathrm{osc}}}
\newcommand{\Scat}{S_{\mathrm{cat}}}
\newcommand{\Spart}{S_{\mathrm{part}}}
\newcommand{\Sfree}{F}

\begin{document}

\title{Operator Trajectories in Neural Partition Space:\\A Formal Language for Bounded Phase-Space Computing}

\author{
Kundai Farai Sachikonye\\
AIMe Registry for Artificial Intelligence\\
\texttt{kundai.sachikonye@bitspark.com}
}

\date{\today}

\maketitle


\begin{abstract}
We develop the \emph{Neural Partition Language and Computing} framework from three axioms: bounded phase space ($\mu(\Omega) < \infty$), no null state (oscillation from categorical necessity), and finite observational resolution ($\delta > 0$). From these axioms we derive partition coordinates $(n, \ell, m, s)$ with shell capacity $C(n) = 2n^2$, and prove the \emph{Triple Equivalence Theorem}: oscillatory, categorical, and partition entropies coincide as $\Sosc = \Scat = \Spart = \kB M \ln n$. The S-entropy coordinate space $(\Sk, \St, \Se) \in [0,1]^3$ provides a navigation metric with ternary trisection yielding $O(\log_3 N)$ search complexity---a 37\% reduction over binary search. We introduce \emph{trajectory-terminus-memory triples} $\Mstate = (\gamma, \Gstate_f, M)$ as complete system descriptors and prove the \emph{Trajectory Non-Identity Theorem}: different paths to the same terminus yield different system states. \emph{Poincar\'{e} Computing} inverts the standard computational paradigm by specifying completion conditions (terminus) and deriving trajectories backward, with applications to drug design and treatment planning. We define the \emph{Pharmacological Neural Partition Language} (pNPL), a typed operator algebra comprising circuit types, pharmacological types, and observable types, with operators for aperture installation, regime classification, coupling modulation, frequency matching, and variance minimization. The free energy $F = \kB T \cdot \sigma^2(\phi)$ defines a therapeutic floor $\sigma^2_{\min} = \kB T / K_{\mathrm{coupling}}$. A five-level frequency hierarchy spanning 18 orders of magnitude connects H$^+$ dynamics ($4.06 \times 10^{13}$ Hz) to behavioral timescales through adiabatic decoupling. Consciousness emerges as the intersection of decay curves $C(t) = \Pdecay \cap \Tdecay$. All 10 validation claims of the pNPL type system pass, including $C(n) = 2n^2$ verification, S-entropy triangle inequality (0 violations in 500 tests), synchronization onset at $K_c = 2\sigma/\pi$, and variance scaling $\sigma^2_{\min} \propto 1/K$ with unit slope in log-log.
\end{abstract}



%% ============================================================
%% SECTION 1: INTRODUCTION
%% ============================================================
\section{Introduction}
\label{sec:introduction}

The central problem of computational neuroscience is this: neural systems are bounded, finite, and dissipative, yet they compute, represent, and predict with extraordinary fidelity~\cite{buzsaki2006rhythms,fries2005mechanism}. Classical computing theory, rooted in unbounded Turing machines, provides an inadequate formal language for such systems. The free-energy principle~\cite{friston2010free} and integrated information theory~\cite{tononi2016integrated} each capture aspects of bounded neural computation, but neither provides a complete formal language with types, operators, and composition rules.

We address this gap by developing the \emph{Neural Partition Language and Computing} (NPLC) framework from first principles. The framework rests on three axioms that any physical computing substrate must satisfy: bounded phase space, the impossibility of a null dynamical state, and finite observational resolution. From these axioms alone, we derive:

\begin{enumerate}[noitemsep]
\item Partition coordinates $(n, \ell, m, s)$ with provable shell capacity $C(n) = 2n^2$.
\item A triple equivalence among oscillatory, categorical, and partition entropies.
\item An S-entropy coordinate space enabling $O(\log_3 N)$ navigation.
\item Trajectory-terminus-memory triples as complete state descriptors.
\item Poincar\'{e} Computing: backward determination from completion conditions.
\item The Pharmacological Neural Partition Language (pNPL) with a complete type system and operator algebra.
\end{enumerate}

The key insight is that boundedness is not a limitation but a \emph{computational resource}. Finite phase space forces discretization (partitions), discretization enables counting (entropy), and counting enables navigation (search). The NPLC framework makes this chain of implications precise and computationally exploitable.

Section~\ref{sec:axioms} establishes the three foundational axioms and derives the triple equivalence. Section~\ref{sec:partition} develops partition coordinates and S-entropy space. Section~\ref{sec:ttm} introduces trajectory-terminus-memory triples. Section~\ref{sec:poincare} presents Poincar\'{e} Computing. Section~\ref{sec:pnpl} defines the pNPL type system and operators. Section~\ref{sec:algebra} establishes composition rules. Section~\ref{sec:frequency} develops the frequency hierarchy. Section~\ref{sec:variance} derives variance minimization results. Section~\ref{sec:validation} presents computational validation. Section~\ref{sec:consciousness} addresses consciousness and time. Sections~\ref{sec:discussion} and~\ref{sec:conclusion} discuss implications and conclude.

%% ============================================================
%% SECTION 2: FOUNDATIONAL AXIOMS AND TRIPLE EQUIVALENCE
%% ============================================================
\section{Foundational Axioms and Triple Equivalence}
\label{sec:axioms}

\subsection{The Three Axioms}

We begin from three axioms that characterize any physical computing substrate.

\begin{axiom}[Bounded Phase Space]
\label{ax:bounded}
The system occupies a phase space $\Omega$ with finite Liouville measure:
\begin{equation}
\mu(\Omega) < \infty.
\label{eq:bounded}
\end{equation}
\end{axiom}

\noindent\textit{Justification.} Every physical system has finite energy and spatial extent~\cite{callen1985thermodynamics,landau1980statistical}. A brain occupies $\sim 1400\,\text{cm}^3$ at temperature $T \approx 310\,\text{K}$, bounding the accessible phase-space volume. This is not an approximation; it is a consequence of finite energy in finite geometry~\cite{jaynes2003probability}.

\begin{axiom}[No Null State]
\label{ax:nonull}
No state in $\Omega$ has zero dynamical activity. Equivalently, the system oscillates from categorical necessity:
\begin{equation}
\forall\, x \in \Omega: \quad \dot{x} \neq 0.
\label{eq:nonull}
\end{equation}
\end{axiom}

\noindent\textit{Justification.} A neural system at absolute zero is no longer a neural system---it is dead~\cite{nicholls2013bioenergetics}. Metabolic maintenance requires continuous H$^+$ flux~\cite{mitchell1961coupling,wikstrom2015proton}, ion channel activity~\cite{doyle1998structure,hille2001ion}, and oscillatory dynamics~\cite{buzsaki2004neuronal}. The zero state is categorically excluded: it belongs to a different ontological class (non-living matter). Oscillation is not a property of the system but a \emph{defining condition} for membership in the category ``neural system.''

\begin{axiom}[Finite Observational Resolution]
\label{ax:resolution}
Every measurement of the system has finite precision $\delta > 0$:
\begin{equation}
\delta > 0.
\label{eq:resolution}
\end{equation}
\end{axiom}

\noindent\textit{Justification.} Physical measurement apparatus has finite resolution~\cite{landauer1961irreversibility,szilard1929entropy}. In neural systems, this is set by ion channel conductance quantization~\cite{doyle1998structure}, synaptic vesicle discreteness, and the thermal noise floor $\kB T$~\cite{callen1985thermodynamics}. Even an ideal observer internal to the system is subject to the resolution set by $\kB T$.

\subsection{Consequences of the Three Axioms}

\begin{theorem}[Forced Partitioning]
\label{thm:partition}
Axioms~\ref{ax:bounded}--\ref{ax:resolution} together imply that $\Omega$ is partitioned into a finite number of distinguishable states $M$:
\begin{equation}
M = \left\lfloor \frac{\mu(\Omega)}{\delta^{d}} \right\rfloor < \infty
\label{eq:M}
\end{equation}
where $d = \dim(\Omega)$.
\end{theorem}

\begin{proof}
By Axiom~\ref{ax:bounded}, $\mu(\Omega) < \infty$. By Axiom~\ref{ax:resolution}, each distinguishable cell has minimum volume $\delta^d > 0$. The maximum number of non-overlapping cells is $M = \lfloor \mu(\Omega)/\delta^d \rfloor$, which is finite. By Axiom~\ref{ax:nonull}, the system must occupy one of these cells at every instant (no cell is dynamically forbidden). Therefore $\Omega$ is partitioned into exactly $M$ distinguishable, dynamically accessible states.
\end{proof}

\begin{theorem}[Oscillatory Entropy]
\label{thm:sosc}
For a system of $M$ coupled oscillators with bounded phase space, the oscillatory entropy (measuring accessible phase configurations) is:
\begin{equation}
\Sosc = \kB \ln W_{\mathrm{osc}}
\label{eq:sosc_def}
\end{equation}
where $W_{\mathrm{osc}}$ is the number of distinguishable oscillatory configurations. For $M$ oscillators each with $n$ accessible phase levels:
\begin{equation}
\Sosc = \kB M \ln n.
\label{eq:sosc}
\end{equation}
\end{theorem}

\begin{proof}
Each of $M$ oscillators can occupy any of $n$ phase-discretized levels (by Theorem~\ref{thm:partition} applied to the phase coordinate). The oscillators are distinguishable (they occupy distinct spatial locations in the neural substrate~\cite{buzsaki2006rhythms}). The total number of configurations is $W_{\mathrm{osc}} = n^M$. Applying the Boltzmann formula~\cite{boltzmann1896lectures,jaynes1957information}: $\Sosc = \kB \ln(n^M) = \kB M \ln n$.
\end{proof}

\begin{theorem}[Categorical Entropy]
\label{thm:scat}
The categorical entropy, measuring the number of distinguishable categories into which the system can be classified, is:
\begin{equation}
\Scat = \kB M \ln n.
\label{eq:scat}
\end{equation}
\end{theorem}

\begin{proof}
By Axiom~\ref{ax:nonull}, each oscillator is categorically active---it belongs to one of $n$ distinguishable categorical states (determined by its phase-amplitude bin). The categorical state of the full system is the product of individual categorical assignments. Since oscillators are distinguishable and each has $n$ categories: $W_{\mathrm{cat}} = n^M$, giving $\Scat = \kB \ln(n^M) = \kB M \ln n$. The key step is that categorical necessity (Axiom~\ref{ax:nonull}) ensures every oscillator contributes exactly one categorical assignment---there are no ``empty'' slots.
\end{proof}

\begin{theorem}[Partition Entropy]
\label{thm:spart}
The partition entropy, counting the number of distinguishable ways to distribute $M$ distinguishable elements across $n$ partition shells, is:
\begin{equation}
\Spart = \kB M \ln n.
\label{eq:spart}
\end{equation}
\end{theorem}

\begin{proof}
Each of $M$ distinguishable elements is assigned to one of $n$ shells. The number of such assignments is $W_{\mathrm{part}} = n^M$ (each element independently chooses from $n$ shells). Therefore $\Spart = \kB \ln(n^M) = \kB M \ln n$.
\end{proof}

We now state the central identity of the framework.

\begin{theorem}[Triple Equivalence]
\label{thm:triple}
For a bounded phase-space system with $M$ oscillators and $n$ accessible levels, the three entropies are identical:
\begin{equation}
\boxed{\Sosc = \Scat = \Spart = \kB M \ln n}
\label{eq:triple}
\end{equation}
\end{theorem}

\begin{proof}
Follows directly from Theorems~\ref{thm:sosc},~\ref{thm:scat}, and~\ref{thm:spart}. The three counting problems---oscillatory configurations, categorical assignments, and partition distributions---are isomorphic. Each reduces to distributing $M$ distinguishable items into $n$ distinguishable bins, yielding $n^M$ microstates and entropy $\kB M \ln n$.
\end{proof}

\begin{remark}
The triple equivalence is not a coincidence but a \emph{theorem} forced by the three axioms. Boundedness forces discretization, discretization forces all three counting problems into the same combinatorial structure, and the no-null-state axiom ensures every element participates. This is the foundational identity upon which the entire NPLC framework rests.
\end{remark}

\subsection{Temperature Factorization}

\begin{theorem}[Temperature Factorization]
\label{thm:temp_factor}
Any observable $O$ in the NPLC framework factorizes as:
\begin{equation}
O = (\kB T) \cdot F(M, n)
\label{eq:temp_factor}
\end{equation}
where $\kB T$ sets the energy scale and $F(M,n)$ is a dimensionless structural function depending only on the partition parameters.
\end{theorem}

\begin{proof}
By the triple equivalence, the entropy $S = \kB M \ln n$ is the fundamental extensive quantity. The conjugate intensive variable is temperature $T$ (from $dU = T\,dS$)~\cite{callen1985thermodynamics}. Any thermodynamic observable derives from the free energy $\mathcal{F} = U - TS$, and $TS = \kB T \cdot M \ln n$. Since all structural information is encoded in $M$ and $n$, any observable decomposes as $O = (\kB T) \cdot F(M,n)$ where $F$ is purely combinatorial.
\end{proof}

%% ============================================================
%% SECTION 3: PARTITION COORDINATES AND S-ENTROPY SPACE
%% ============================================================
\section{Partition Coordinates and S-Entropy Space}
\label{sec:partition}

\subsection{Partition Coordinates $(n, \ell, m, s)$}

\begin{definition}[Partition Coordinates]
\label{def:coords}
A partition state is specified by four quantum numbers:
\begin{itemize}[noitemsep]
\item $n \in \{1, 2, 3, \ldots\}$: principal partition number (shell)
\item $\ell \in \{0, 1, \ldots, n-1\}$: angular partition number (sub-shell)
\item $m \in \{-\ell, -\ell+1, \ldots, \ell\}$: magnetic partition number (orientation)
\item $s \in \{-\tfrac{1}{2}, +\tfrac{1}{2}\}$: spin partition number (binary state)
\end{itemize}
\end{definition}

\begin{theorem}[Shell Capacity]
\label{thm:capacity}
The capacity of shell $n$ is:
\begin{equation}
\boxed{C(n) = 2n^2}
\label{eq:capacity}
\end{equation}
\end{theorem}

\begin{proof}
We count the number of valid $(n, \ell, m, s)$ tuples for fixed $n$.

\textbf{Step 1.} For each $\ell$, the magnetic number $m$ ranges from $-\ell$ to $+\ell$, giving $2\ell + 1$ values.

\textbf{Step 2.} For each $(n, \ell, m)$, the spin $s$ takes 2 values.

\textbf{Step 3.} Sum over all allowed $\ell$:
\begin{align}
C(n) &= \sum_{\ell=0}^{n-1} 2(2\ell + 1) \notag \\
&= 2 \sum_{\ell=0}^{n-1} (2\ell + 1) \notag \\
&= 2 \left[ \sum_{\ell=0}^{n-1} 2\ell + \sum_{\ell=0}^{n-1} 1 \right] \notag \\
&= 2 \left[ 2 \cdot \frac{(n-1)n}{2} + n \right] \notag \\
&= 2 \left[ n(n-1) + n \right] \notag \\
&= 2 \left[ n^2 - n + n \right] \notag \\
&= 2n^2. \qedhere
\end{align}
\end{proof}

\begin{corollary}[Cumulative Capacity]
\label{cor:cumcap}
The total number of states up to and including shell $n$ is:
\begin{equation}
N(n) = \sum_{k=1}^{n} 2k^2 = \frac{n(n+1)(2n+1)}{3}.
\label{eq:cumcap}
\end{equation}
\end{corollary}

\begin{proof}
Using the identity $\sum_{k=1}^n k^2 = n(n+1)(2n+1)/6$:
$N(n) = 2 \sum_{k=1}^n k^2 = 2 \cdot \frac{n(n+1)(2n+1)}{6} = \frac{n(n+1)(2n+1)}{3}$.
\end{proof}

\begin{figure*}[!htbp]
    \centering
    \includegraphics[width=0.9\textwidth]{figures/panel1_partition_coordinate_space.png}
    \caption{\textbf{Partition Coordinate Space.}
    (\textbf{A}) Three-dimensional scatter of all valid partition states $(n, \ell, m)$ for principal quantum numbers $n = 1$--$5$, with each spin degeneracy ($s = \pm\tfrac{1}{2}$) plotted separately. Point color encodes $n$ and size scales with $|m| + 1$. The discrete lattice structure visualizes the $C(n) = 2n^2$ capacity formula: the $n = 5$ shell contains $50$ states spanning $\ell = 0$--$4$ and $m = -4$--$+4$.
    (\textbf{B}) Partition capacity $C(n) = 2n^2$ as bar heights for $n = 1$--$10$, with the continuous $2n^2$ curve overlaid (dashed). Orange markers highlight the seven levels validated against explicit state enumeration, where $C_\mathrm{formula} = C_\mathrm{enumerated}$ for all $n$.
    (\textbf{C}) Three-dimensional surface of the partition depth $M(K, k) = K \log_3 k$ as a function of the number of hierarchical levels $K$ and branching factor $k$. The surface demonstrates that depth grows logarithmically in branching but linearly in levels, establishing the resource scaling for the type system.
    (\textbf{D}) State density heatmap for the $n = 5$ shell in the $(\ell, m)$ plane. Each occupied cell has degeneracy $2$ (spin), with the triangular occupancy pattern $|m| \leq \ell$ clearly visible. Unoccupied cells ($\ell < |m|$) appear as zero, confirming the angular momentum selection rule.}
    \label{fig:partition_coordinate_space}
\end{figure*}

\subsection{S-Entropy Coordinate Space}

\begin{definition}[S-Entropy Coordinates]
\label{def:sentropy}
The S-entropy coordinate space is $\Sspace = [0,1]^3$ with coordinates:
\begin{itemize}[noitemsep]
\item $\Sk \in [0,1]$: \emph{knowledge entropy}---uncertainty in state identification
\item $\St \in [0,1]$: \emph{temporal entropy}---uncertainty in temporal sequencing
\item $\Se \in [0,1]$: \emph{evolution entropy}---uncertainty in trajectory through partition space
\end{itemize}
A system state maps to $\Scoord = (\Sk, \St, \Se) \in [0,1]^3$.
\end{definition}

\begin{definition}[S-Distance]
\label{def:sdistance}
The distance between two states in S-entropy space is:
\begin{equation}
\dS(\Scoord_1, \Scoord_2) = \sqrt{(\Delta \Sk)^2 + (\Delta \St)^2 + (\Delta \Se)^2}
\label{eq:sdist}
\end{equation}
where $\Delta \Sk = \Sk^{(2)} - \Sk^{(1)}$, etc.
\end{definition}

\begin{proposition}[S-Distance Bound]
\label{prop:sbound}
For all states $\Scoord_1, \Scoord_2 \in [0,1]^3$:
\begin{equation}
0 \leq \dS(\Scoord_1, \Scoord_2) \leq \sqrt{3}.
\label{eq:sbound}
\end{equation}
\end{proposition}

\begin{proof}
The minimum is achieved when $\Scoord_1 = \Scoord_2$. The maximum is achieved when one state is at $(0,0,0)$ and the other at $(1,1,1)$: $\dS = \sqrt{1+1+1} = \sqrt{3}$.
\end{proof}

\subsection{Ternary Trisection Navigation}

\begin{theorem}[Trisection Search Complexity]
\label{thm:trisection}
Navigation in the S-entropy cube $[0,1]^3$ via ternary trisection achieves:
\begin{equation}
\text{Search complexity} = O(\log_3 N)
\label{eq:trisection}
\end{equation}
representing a 37\% reduction in comparisons versus binary search.
\end{theorem}

\begin{proof}
At each step, trisection divides each coordinate into three equal intervals, reducing the search volume by a factor of $3^3 = 27$ in three dimensions (or by 3 in one dimension). After $k$ steps, the remaining fraction is $(1/3)^k$. To resolve $N$ states requires $(1/3)^k \leq 1/N$, giving $k \geq \log_3 N$.

The reduction factor relative to binary search is:
\begin{equation}
\frac{\log_3 N}{\log_2 N} = \frac{\ln 2}{\ln 3} = \frac{1}{\log_2 3} \approx 0.631.
\end{equation}
This is a $(1 - 0.631) \times 100\% \approx 37\%$ reduction.
\end{proof}

\begin{remark}
Ternary trisection is natural for the three-coordinate S-entropy space: each query eliminates two-thirds of the search volume along one axis. This is exploited in the pNPL regime classification operators (Section~\ref{sec:pnpl}).
\end{remark}

%% ============================================================
%% SECTION 4: TRAJECTORY-TERMINUS-MEMORY TRIPLES
%% ============================================================
\section{Trajectory-Terminus-Memory Triples}
\label{sec:ttm}

\subsection{Complete State Description}

\begin{definition}[Trajectory-Terminus-Memory Triple]
\label{def:ttm}
The complete state of a neural partition system is the triple:
\begin{equation}
\Mstate = (\gamma, \Gstate_f, M)
\label{eq:ttm}
\end{equation}
where:
\begin{itemize}[noitemsep]
\item $\gamma: [0,T] \to \Sspace$ is the \emph{trajectory} through S-entropy space
\item $\Gstate_f = \gamma(T)$ is the \emph{terminus} (current configuration)
\item $M = \int_0^T \dot{\Hfield}(t)\,dt$ is the \emph{memory} (accumulated H$^+$ field change)
\end{itemize}
\end{definition}

The terminus $\Gstate_f$ alone is insufficient because it discards the history that produced it. The memory $M$ encodes the cumulative emotional-metabolic context~\cite{damasio1994descartes,mitchell1961coupling}. Two systems at the same terminus with different memories are physically different---their H$^+$ field configurations differ, which influences subsequent dynamics.

\begin{theorem}[Trajectory Non-Identity]
\label{thm:nonidentity}
Let $\Mstate_1 = (\gamma_1, \Gstate_f, M_1)$ and $\Mstate_2 = (\gamma_2, \Gstate_f, M_2)$ be two states with the same terminus $\Gstate_f$ but different trajectories $\gamma_1 \neq \gamma_2$. Then $\Mstate_1 \neq \Mstate_2$.
\end{theorem}

\begin{proof}
Since $\gamma_1 \neq \gamma_2$, the accumulated H$^+$ field changes differ:
\begin{equation}
M_1 = \int_0^{T_1} \dot{\Hfield}[\gamma_1(t)]\,dt \neq \int_0^{T_2} \dot{\Hfield}[\gamma_2(t)]\,dt = M_2
\end{equation}
because $\dot{\Hfield}$ depends on the trajectory (the H$^+$ production rate varies with the sequence of partition states traversed~\cite{wikstrom2015proton,nicholls2013bioenergetics}). Different trajectories sample different metabolic loads, ion channel activation sequences~\cite{hille2001ion}, and protein conformational histories~\cite{hartl2011molecular}. Since $M_1 \neq M_2$, we have $\Mstate_1 \neq \Mstate_2$ even though $\Gstate_{f,1} = \Gstate_{f,2}$.
\end{proof}

\begin{corollary}[Path Dependence of Neural States]
\label{cor:pathdep}
The current state of a neural partition system cannot be determined from instantaneous measurements alone. The trajectory history is a necessary component of the state description.
\end{corollary}

\begin{remark}
This is the mathematical formalization of a deep intuition: how you arrived at a thought matters. The same conclusion reached through careful reasoning versus snap judgment constitutes a different mental state, with different subsequent dynamics~\cite{tulving1972episodic,eichenbaum2000hippocampus}. The trajectory-terminus-memory triple makes this precise.
\end{remark}

\subsection{Trajectory Space Structure}

\begin{definition}[Trajectory Equivalence]
\label{def:traj_equiv}
Two trajectories $\gamma_1, \gamma_2$ are \emph{M-equivalent} if they produce the same memory:
\begin{equation}
\gamma_1 \sim_M \gamma_2 \iff M[\gamma_1] = M[\gamma_2].
\end{equation}
The equivalence class $[\gamma]_M$ is the set of all trajectories producing the same accumulated H$^+$ change.
\end{definition}

\begin{proposition}[Trajectory Classes are Measure-Zero]
\label{prop:measure_zero}
For generic $\dot{\Hfield}$, the M-equivalence class $[\gamma]_M$ has measure zero in the space of all trajectories connecting the same endpoints.
\end{proposition}

\begin{proof}
The memory functional $M[\gamma] = \int_0^T \dot{\Hfield}[\gamma(t)]\,dt$ is a real-valued functional on the infinite-dimensional space of paths. For generic (non-degenerate) $\dot{\Hfield}$, the level sets $M^{-1}(c)$ are codimension-1 submanifolds, hence measure zero.
\end{proof}

\begin{figure*}[!htbp]
    \centering
    \includegraphics[width=0.9\textwidth]{figures/panel3_trajectory_terminus_memory.png}
    \caption{\textbf{Trajectory-Terminus-Memory Triples.}
    (\textbf{A}) Two three-dimensional trajectories $\gamma_1$ (blue) and $\gamma_2$ (red) converging to a common terminus $\Gamma_f$ (gold star) in the $(S_k, S_t, S_e)$ space. Despite originating from different pathological states, both trajectories arrive at the same therapeutic fixed point, demonstrating the backward-determination principle: the terminus constrains the trajectory, not vice versa.
    (\textbf{B}) Accumulated memory $M(t) = \int_0^t \|\dot{\gamma}(\tau)\| \, d\tau$ for both trajectories. Trajectory $\gamma_2$ accumulates more total memory (longer path length) due to its more circuitous route through S-entropy space, while both converge to finite values at $t = 1$. Horizontal dashed lines mark the final memory $M(1)$ for each path.
    (\textbf{C}) Scatter of $50$ simulated trajectory endpoints in the $(S_k, S_e)$ plane, colored by accumulated memory $M$. The clustering near the terminus region $(0.25$--$0.55, 0.2$--$0.5)$ with variable memory values confirms that distinct therapeutic histories can converge to similar outcomes with different memory footprints.
    (\textbf{D}) Poincar\'{e} near-return distance $\|\gamma(t) - \gamma(0)\|$ over extended time ($t = 0$--$50$). The quasi-periodic oscillations with secular drift demonstrate that the trajectory does not return exactly to its origin, consistent with the irreversibility of memory accumulation in the partition framework.}
    \label{fig:trajectory_terminus_memory}
\end{figure*}


%% ============================================================
%% SECTION 5: POINCARE COMPUTING
%% ============================================================
\section{Poincar\'{e} Computing: Backward Determination}
\label{sec:poincare}

\subsection{The Backward Paradigm}

Standard computation proceeds forward: given initial conditions, evolve the system and observe the output. Poincar\'{e} Computing~\cite{poincare1890probleme,poincare1892methodes} inverts this paradigm: specify the desired \emph{completion condition} (terminus) and determine the trajectory that achieves it.

\begin{definition}[Poincar\'{e} Computing Problem]
\label{def:poincare}
Given a terminus $\Gstate_f \in \Sspace$ and a constraint set $\mathcal{C}$, the Poincar\'{e} Computing Problem is to find the trajectory class:
\begin{equation}
\exists!\,[\gamma]: \quad \gamma(T) = \Gstate_f \;\wedge\; \text{Constraints}(\gamma) = \text{true}
\label{eq:poincare}
\end{equation}
where uniqueness is up to M-equivalence.
\end{definition}

\begin{theorem}[Existence of Backward Solutions]
\label{thm:backward}
Under the three axioms, the Poincar\'{e} Computing Problem has a solution whenever $\Gstate_f$ lies in the accessible region of $\Sspace$ and the constraints are compatible with the system dynamics.
\end{theorem}

\begin{proof}
By Axiom~\ref{ax:bounded}, the accessible region is compact. By Axiom~\ref{ax:nonull}, the dynamics are everywhere non-degenerate, so the flow map $\Phi_t: \Omega \to \Omega$ is a diffeomorphism for each $t$. By Axiom~\ref{ax:resolution}, the target $\Gstate_f$ has a resolution neighborhood of volume $\delta^d > 0$. We must show that the pre-image $\Phi_{-T}(B_\delta(\Gstate_f))$ is non-empty for some $T > 0$.

Since $\Phi_t$ preserves the Liouville measure (or, for dissipative systems, contracts it), and the accessible region is compact with positive measure, the Poincar\'{e} recurrence theorem~\cite{poincare1890probleme,arnold1989mathematical} guarantees that almost every trajectory returns arbitrarily close to any accessible state. Therefore, for any $\Gstate_f$ in the accessible region, there exists a trajectory reaching $B_\delta(\Gstate_f)$.

The constraint compatibility condition excludes pathological cases (e.g., requiring the trajectory to simultaneously occupy two disjoint regions).
\end{proof}

\subsection{Application: Drug Design}

In Poincar\'{e} drug design, one specifies the therapeutic terminus---the desired S-entropy configuration representing a healthy brain state---and computes backward to determine the drug-induced trajectory modifications required.

\begin{definition}[Therapeutic Terminus]
\label{def:therapeutic_terminus}
A therapeutic terminus is a target state $\Gstate_f^* \in \Sspace$ satisfying:
\begin{enumerate}[noitemsep]
\item $\Gstate_f^*$ lies in the healthy regime of S-entropy space
\item $\Gstate_f^*$ is dynamically stable (attracting)
\item The trajectory from the pathological state to $\Gstate_f^*$ satisfies safety constraints
\end{enumerate}
\end{definition}

\begin{definition}[Treatment Planning Problem]
\label{def:treatment}
Given a pathological initial state $\Mstate_0 = (\gamma_0, \Gstate_0, M_0)$ and a therapeutic terminus $\Gstate_f^*$, find a drug operator $\Pop_{\mathrm{drug}}$ such that:
\begin{equation}
\Pop_{\mathrm{drug}} \circ \Mstate_0 = (\gamma^*, \Gstate_f^*, M^*)
\end{equation}
where $\gamma^*$ satisfies the safety constraints and $M^*$ is within tolerable memory bounds.
\end{definition}

%% ============================================================
%% SECTION 6: THE PHARMACOLOGICAL NEURAL PARTITION LANGUAGE
%% ============================================================
\section{The Pharmacological Neural Partition Language}
\label{sec:pnpl}

We now define pNPL as a typed formal language for expressing pharmacological computations on neural partition systems.

\subsection{Type System}

\subsubsection{Circuit Types}

\begin{definition}[Circuit Types]
\label{def:circuit_types}
The circuit type system comprises:
\begin{align}
&\texttt{PartitionCoord} : (n, \ell, m, s) \text{ with } C(n) = 2n^2 \label{eq:type_pc} \\
&\texttt{SCoord} : (\Sk, \St, \Se) \in [0,1]^3 \label{eq:type_sc} \\
&\texttt{Regime} : \{\texttt{coherent}, \texttt{turbulent}, \notag \\
&\qquad\qquad\quad \texttt{cascade}, \texttt{aperture}, \texttt{sync}\} \label{eq:type_regime} \\
&\texttt{StructuralFactor} : \mathcal{S} \in \mathbb{R}^+ \label{eq:type_sf}
\end{align}
\end{definition}

\begin{definition}[Regime Classification]
\label{def:regime_class}
The five regimes correspond to ranges of the Kuramoto order parameter $R \in [0,1]$~\cite{kuramoto1984chemical,strogatz2000kuramoto}:
\begin{align}
&\texttt{turbulent}: && R \in [0, 0.2) \notag \\
&\texttt{cascade}: && R \in [0.2, 0.4) \notag \\
&\texttt{aperture}: && R \in [0.4, 0.6) \notag \\
&\texttt{coherent}: && R \in [0.6, 0.8) \notag \\
&\texttt{sync}: && R \in [0.8, 1.0] \notag
\end{align}
This classification is exhaustive (covers all $R \in [0,1]$) and monotonic (increasing $R$ maps to regimes of increasing order).
\end{definition}

\subsubsection{Pharmacological Types}

\begin{definition}[Pharmacological Types]
\label{def:pharm_types}
\begin{align}
&\texttt{Aperture} : (A, \textit{type}, \alpha) \notag \\
&\quad \text{where } \textit{type} \in \{\texttt{monopole}, \texttt{dipole}, \texttt{quadrupole}\} \notag \\
&\quad \text{and } \alpha \in \mathbb{R} \text{ (aperture parameter)} \label{eq:type_ap} \\[4pt]
&\texttt{CouplingMod} : \Delta K \in \mathbb{R} \label{eq:type_cm} \\[4pt]
&\texttt{DrugTrajectory} : \Phi_{\mathrm{drug}} : \Mstate \to \Mstate \label{eq:type_dt} \\[4pt]
&\texttt{TransitionPath} : [\texttt{Regime}] \text{ (ordered sequence)} \label{eq:type_tp}
\end{align}
\end{definition}

The aperture types classify drug-receptor interactions by their symmetry: monopole (single-site agonism), dipole (allosteric modulation with two-site geometry), and quadrupole (complex multi-site interactions)~\cite{kenakin2017pharmacology,roth2004multiplicity}.

\subsubsection{Observable Types}

\begin{definition}[Observable Types]
\label{def:obs_types}
\begin{align}
&\texttt{OrderParam} : R \in [0,1] \label{eq:type_op} \\
&\texttt{PhaseVariance} : \sigma^2 \in \mathbb{R}^+ \label{eq:type_pv} \\
&\texttt{PLV} : \text{phase-locking value} \in [0,1] \label{eq:type_plv} \\
&\texttt{OnsetTime} : \tau_{\mathrm{onset}} \in \mathbb{R}^+ \label{eq:type_ot}
\end{align}
\end{definition}

The phase-locking value (PLV) quantifies the consistency of phase relationships between oscillator pairs~\cite{lachaux1999measuring}.

\subsection{Core Operators}

\subsubsection{Aperture Operations}

\begin{definition}[Aperture Operators]
\label{def:aperture_ops}
\begin{align}
&\textsc{Install\_Aperture} : \texttt{Aperture} \times \texttt{SCoord} \to \texttt{SCoord} \notag \\
&\textsc{Classify\_Aperture} : \texttt{Aperture} \to \textit{type} \notag \\
&\textsc{Catalytic\_Factor} : \texttt{Aperture} \to \texttt{StructuralFactor} \notag
\end{align}
$\textsc{Install\_Aperture}(A, \Scoord)$ modifies the S-entropy coordinates by inserting a drug-induced aperture at position $\Scoord$. $\textsc{Classify\_Aperture}(A)$ returns the symmetry type. $\textsc{Catalytic\_Factor}(A)$ returns the structural factor $\mathcal{S} \in (0,1]$ by which the aperture modifies transition rates.
\end{definition}

\subsubsection{Regime Operations}

\begin{definition}[Regime Operators]
\label{def:regime_ops}
\begin{align}
&\textsc{Regime\_Classify} : \texttt{OrderParam} \to \texttt{Regime} \notag \\
&\textsc{Structural\_Factor} : \texttt{Regime} \to \texttt{StructuralFactor} \notag \\
&\textsc{Delta\_S} : \texttt{Regime} \times \texttt{Regime} \to \mathbb{R} \notag
\end{align}
$\textsc{Regime\_Classify}(R)$ maps the order parameter to the appropriate regime. $\textsc{Structural\_Factor}(r)$ returns the structural modification factor for regime $r$. $\textsc{Delta\_S}(r_1, r_2)$ computes the entropy change for a regime transition.
\end{definition}

\subsubsection{Transition Operations}

\begin{definition}[Transition Operators]
\label{def:trans_ops}
\begin{align}
&\textsc{Transition\_Path} : \texttt{Regime} \times \texttt{Regime} \to \texttt{TransitionPath} \notag \\
&\textsc{Onset\_Time} : \texttt{TransitionPath} \to \texttt{OnsetTime} \notag \\
&\textsc{Cross\_Modal\_Test} : \texttt{TransitionPath} \to \{0, 1\} \notag
\end{align}
$\textsc{Transition\_Path}(r_1, r_2)$ computes the ordered sequence of regimes traversed. $\textsc{Onset\_Time}(\pi)$ returns the time for the first regime change. $\textsc{Cross\_Modal\_Test}(\pi)$ returns 1 if the transition path crosses a modal boundary~\cite{varela2001brainweb,engel2001dynamic}.
\end{definition}

\subsubsection{Coupling Operations}

\begin{definition}[Coupling Operators]
\label{def:coupling_ops}
\begin{align}
&\textsc{Modulate\_K} : \texttt{CouplingMod} \times \texttt{SCoord} \to \texttt{SCoord} \notag \\
&\textsc{Critical\_K} : \texttt{SCoord} \to \mathbb{R}^+ \notag \\
&\textsc{Phase\_Lock} : \texttt{CouplingMod} \to \texttt{PLV} \notag
\end{align}
$\textsc{Modulate\_K}(\Delta K, \Scoord)$ modifies the coupling strength and returns the new S-entropy coordinates. $\textsc{Critical\_K}(\Scoord)$ returns the critical coupling for synchronization onset at the current state~\cite{acebron2005kuramoto}. $\textsc{Phase\_Lock}(\Delta K)$ returns the resulting PLV.
\end{definition}

\subsubsection{Frequency Operations}

\begin{definition}[Frequency Operators]
\label{def:freq_ops}
\begin{align}
&\textsc{Freq\_Match} : \omega_{\mathrm{drug}} \times \omega_{\mathrm{circuit}} \to [0,1] \notag \\
&\textsc{Gear\_Activate} : \omega_{\mathrm{drug}} \times G \to \omega_{\mathrm{therapeutic}} \notag
\end{align}
$\textsc{Freq\_Match}(\omega_d, \omega_c)$ returns the frequency overlap between drug kinetics and circuit dynamics. $\textsc{Gear\_Activate}(\omega_d, G)$ activates a gear network with amplification factor $G$, producing $\omega_{\mathrm{therapeutic}} = G \times \omega_{\mathrm{drug}}$.
\end{definition}

\subsubsection{Trajectory Operations}

\begin{definition}[Trajectory Operators]
\label{def:traj_ops}
\begin{align}
&\textsc{Drug\_Modify} : \texttt{DrugTrajectory} \times \Mstate \to \Mstate \notag \\
&\textsc{Variance\_Minimize} : \Mstate \times \sigma^2_{\mathrm{target}} \to \texttt{DrugTrajectory} \notag \\
&\textsc{Complete} : \Mstate \times \Gstate_f \to \texttt{TransitionPath} \notag
\end{align}
$\textsc{Drug\_Modify}(\Phi, \Mstate)$ applies the drug trajectory operator. $\textsc{Variance\_Minimize}(\Mstate, \sigma^2_t)$ finds the drug trajectory that minimizes variance to the target. $\textsc{Complete}(\Mstate, \Gstate_f)$ computes the Poincar\'{e} backward trajectory to terminus $\Gstate_f$.
\end{definition}

\subsubsection{Measurement Operations}

\begin{definition}[Measurement Operators]
\label{def:meas_ops}
\begin{align}
&\textsc{PLV\_Measure} : \Mstate \to \texttt{PLV} \notag \\
&\textsc{Regime\_Map} : \Mstate \to \texttt{Regime}^{N_{\mathrm{regions}}} \notag \\
&\textsc{Categorical\_Therm} : \Mstate \to (\kB T, \mathcal{S}) \notag
\end{align}
$\textsc{PLV\_Measure}(\Mstate)$ extracts the phase-locking value. $\textsc{Regime\_Map}(\Mstate)$ returns the regime classification across all brain regions. $\textsc{Categorical\_Therm}(\Mstate)$ returns the categorical temperature and structural factor.
\end{definition}

\subsection{Algorithm: Antidepressant Treatment in pNPL}

\begin{algorithm}[t]
\caption{Antidepressant Treatment Protocol in pNPL}
\label{alg:antidepressant}
\begin{algorithmic}[1]
\Require Patient state $\Mstate_0 = (\gamma_0, \Gstate_0, M_0)$
\Require Drug parameters: $\Delta K_{\mathrm{SSRI}}$, $\omega_{\mathrm{drug}}$
\Ensure Treated state $\Mstate_f$ in target regime
\State $R_0 \gets \textsc{PLV\_Measure}(\Mstate_0)$
\State $r_0 \gets \textsc{Regime\_Classify}(R_0)$
\State $K_c \gets \textsc{Critical\_K}(\Scoord_0)$
\If{$r_0 = \texttt{turbulent}$}
    \State $\triangleright$ \textit{Phase 1: Establish minimal coherence}
    \State $A_{\mathrm{ssri}} \gets (\Delta K_{\mathrm{SSRI}}, \texttt{monopole}, 0.3)$
    \State $\Scoord_1 \gets \textsc{Install\_Aperture}(A_{\mathrm{ssri}}, \Scoord_0)$
    \State $\Scoord_2 \gets \textsc{Modulate\_K}(\Delta K_{\mathrm{SSRI}}, \Scoord_1)$
    \State $\tau \gets \textsc{Onset\_Time}(\textsc{Transition\_Path}(r_0, \texttt{cascade}))$
    \State \textbf{wait} $\tau$ \Comment{Therapeutic latency}
\EndIf
\State $\triangleright$ \textit{Phase 2: Drive toward coherent regime}
\State $\pi^* \gets \textsc{Transition\_Path}(r_0, \texttt{coherent})$
\ForAll{$(r_i, r_{i+1}) \in \pi^*$}
    \State $\Delta K_i \gets \textsc{Critical\_K}(\Scoord_i) - K_{\mathrm{current}}$
    \State $\Scoord_{i+1} \gets \textsc{Modulate\_K}(\Delta K_i, \Scoord_i)$
    \State $\sigma^2_i \gets \textsc{PLV\_Measure}(\Mstate_i)$
    \If{$\sigma^2_i > \sigma^2_{\mathrm{safe}}$}
        \State $\Phi_{\mathrm{adj}} \gets \textsc{Variance\_Minimize}(\Mstate_i, \sigma^2_{\mathrm{safe}})$
        \State $\Mstate_i \gets \textsc{Drug\_Modify}(\Phi_{\mathrm{adj}}, \Mstate_i)$
    \EndIf
\EndFor
\State $\triangleright$ \textit{Phase 3: Verify and stabilize}
\State $R_f \gets \textsc{PLV\_Measure}(\Mstate_f)$
\State $r_f \gets \textsc{Regime\_Classify}(R_f)$
\State \textbf{assert} $r_f \in \{\texttt{coherent}, \texttt{sync}\}$
\State \Return $\Mstate_f$
\end{algorithmic}
\end{algorithm}

\subsection{Algorithm: Drug Design via Poincar\'{e} Computing}

\begin{algorithm}[t]
\caption{Drug Design via Poincar\'{e} Computing}
\label{alg:poincare_drug}
\begin{algorithmic}[1]
\Require Therapeutic terminus $\Gstate_f^*$, safety constraints $\mathcal{C}$
\Require Pathological population $\{\Mstate_0^{(i)}\}_{i=1}^{N_p}$
\Ensure Drug operator $\Phi_{\mathrm{drug}}^*$
\State $\triangleright$ \textit{Step 1: Backward trajectory computation}
\ForAll{$\Mstate_0^{(i)} \in \{\Mstate_0^{(i)}\}$}
    \State $\pi_i \gets \textsc{Complete}(\Mstate_0^{(i)}, \Gstate_f^*)$
    \State $\tau_i \gets \textsc{Onset\_Time}(\pi_i)$
    \State \textbf{assert} $\textsc{Cross\_Modal\_Test}(\pi_i) = 0$ \Comment{Safety}
\EndFor
\State $\triangleright$ \textit{Step 2: Extract common trajectory structure}
\State $\bar{\pi} \gets \textsc{Mode}(\{\pi_i\})$ \Comment{Most common path}
\State $\bar{\tau} \gets \textsc{Median}(\{\tau_i\})$
\State $\triangleright$ \textit{Step 3: Determine required coupling modifications}
\ForAll{$(r_j, r_{j+1}) \in \bar{\pi}$}
    \State $\Delta K_j \gets \textsc{Delta\_S}(r_j, r_{j+1}) / (\kB T)$
    \State $\mathcal{S}_j \gets \textsc{Structural\_Factor}(r_j)$
    \State $\omega_j \gets \textsc{Freq\_Match}(\omega_{\mathrm{mol}}, \omega_{r_j})$
\EndFor
\State $\triangleright$ \textit{Step 4: Synthesize drug operator}
\State $\Phi_{\mathrm{drug}}^* \gets \prod_j \textsc{Modulate\_K}(\mathcal{S}_j \cdot \Delta K_j, \cdot)$
\State $\triangleright$ \textit{Step 5: Validate}
\ForAll{$\Mstate_0^{(i)}$}
    \State $\Mstate_f^{(i)} \gets \textsc{Drug\_Modify}(\Phi_{\mathrm{drug}}^*, \Mstate_0^{(i)})$
    \State $R_f \gets \textsc{PLV\_Measure}(\Mstate_f^{(i)})$
    \State \textbf{assert} $\dS(\Gstate_f^{(i)}, \Gstate_f^*) < \epsilon$
\EndFor
\State \Return $\Phi_{\mathrm{drug}}^*$
\end{algorithmic}
\end{algorithm}

%% ============================================================
%% SECTION 7: OPERATOR ALGEBRA AND COMPOSITION RULES
%% ============================================================
\section{Operator Algebra and Composition Rules}
\label{sec:algebra}

\subsection{Composition Structure}

The pNPL operators form an algebra under sequential composition.

\begin{definition}[Operator Composition]
\label{def:composition}
For operators $\mathcal{O}_1: A \to B$ and $\mathcal{O}_2: B \to C$, the composition $\mathcal{O}_2 \circ \mathcal{O}_1: A \to C$ is defined by:
\begin{equation}
(\mathcal{O}_2 \circ \mathcal{O}_1)(x) = \mathcal{O}_2(\mathcal{O}_1(x)).
\end{equation}
\end{definition}

\begin{theorem}[Type Safety of Composition]
\label{thm:type_safe}
The composition $\textsc{Install\_Aperture} \circ \textsc{Regime\_Classify} \circ \textsc{Modulate\_K}$ produces valid trajectories, i.e., the output lies in the S-entropy cube $[0,1]^3$.
\end{theorem}

\begin{proof}
$\textsc{Modulate\_K}: \mathbb{R} \times [0,1]^3 \to [0,1]^3$ (coupling modification preserves the S-entropy bounds because the Kuramoto dynamics are bounded~\cite{acebron2005kuramoto}). $\textsc{Regime\_Classify}: [0,1] \to \texttt{Regime}$ (the order parameter $R$ derived from any $\Scoord \in [0,1]^3$ satisfies $R \in [0,1]$ by definition). $\textsc{Install\_Aperture}: \texttt{Aperture} \times [0,1]^3 \to [0,1]^3$ (aperture installation modifies S-entropy coordinates within bounds because the aperture parameter $\alpha$ is normalized). Composing: the input $(\Delta K, \Scoord)$ produces $\Scoord' \in [0,1]^3$, which yields $R \in [0,1]$, which classifies to a valid regime, and the aperture installation on $\Scoord'$ remains in $[0,1]^3$.
\end{proof}

\begin{proposition}[Non-Commutativity]
\label{prop:noncommute}
In general, pNPL operators do not commute:
\begin{equation}
\textsc{Modulate\_K} \circ \textsc{Install\_Aperture} \neq \textsc{Install\_Aperture} \circ \textsc{Modulate\_K}.
\end{equation}
\end{proposition}

\begin{proof}
The aperture modifies the structural factor, which changes the effective coupling landscape. Modulating coupling first on the unmodified landscape, then installing the aperture, produces a different S-entropy trajectory than installing the aperture first (which changes the landscape) and then modulating coupling on the modified landscape. This is physically natural: administering a drug before or after changing the coupling environment yields different outcomes~\cite{stahl2013psychopharmacology}.
\end{proof}

\subsection{Idempotency and Fixed Points}

\begin{proposition}[Regime Classification Idempotency]
\label{prop:idempotent}
$\textsc{Regime\_Classify}$ is idempotent in the sense that applying it to an already-classified state returns the same regime:
\begin{equation}
\textsc{Regime\_Classify}(\textsc{Regime\_Classify}(R)) = \textsc{Regime\_Classify}(R).
\end{equation}
\end{proposition}

\begin{proposition}[Synchronization Fixed Point]
\label{prop:sync_fix}
For $K > K_c$, the \texttt{sync} regime is a fixed point of $\textsc{Modulate\_K}(\Delta K, \cdot)$ for $\Delta K > 0$:
\begin{equation}
\textsc{Regime\_Classify}(\textsc{Modulate\_K}(\Delta K, \Scoord_{\mathrm{sync}})) = \texttt{sync}.
\end{equation}
\end{proposition}

\begin{proof}
In the \texttt{sync} regime, $R \geq 0.8$. Increasing coupling ($\Delta K > 0$) can only increase or maintain the order parameter (by the monotonicity of the Kuramoto order parameter in coupling strength~\cite{strogatz2000kuramoto}). Therefore $R' \geq R \geq 0.8$, and the regime remains \texttt{sync}.
\end{proof}

\begin{figure*}[!htbp]
    \centering
    \includegraphics[width=0.9\textwidth]{figures/panel4_operator_composition_drug.png}
    \caption{\textbf{Operator Composition and Drug Action.}
    (\textbf{A}) Three-dimensional trajectory in (step, $R$, $\sigma^2$) space for a $200$-step simulation with drug application at step $100$ (red triangle). Color encodes regime classification: turbulent (red), aperture-dominated (orange), cascade (green), coherent (cyan), phase-locked (blue). The trajectory shows the system ascending from the turbulent manifold ($R \approx 0.13$, high variance) through successive regimes to phase-locked coherence ($R = 0.999$).
    (\textbf{B}) Order parameter $R(t)$ over $200$ operator steps with regime-colored background bands. The vertical dashed line at step $100$ marks the drug operator application (APERTURE $\circ$ REGIME $\circ$ COUPLE). The pre-drug turbulent baseline ($R = 0.135$) transitions through all five regimes, reaching $R = 0.999$ post-drug, validating the operator composition sequence.
    (\textbf{C}) Heatmap of the post-drug order parameter $R_\mathrm{post}$ as a function of pre-drug state $R_\mathrm{pre}$ and structural factor modification $\Delta S$. The color gradient from low (dark) to high (bright) $R_\mathrm{post}$ shows that larger $\Delta S$ corrections are required to achieve coherence from deeper turbulent states.
    (\textbf{D}) Bar chart of regime distribution across $1001$ uniformly sampled $R$ values: turbulent ($301$), aperture-dominated ($200$), cascade ($299$), coherent ($150$), phase-locked ($51$). The monotonically decreasing count toward higher regimes reflects the progressive narrowing of each regime's $R$-interval.}
    \label{fig:operator_composition}
\end{figure*}

%% ============================================================
%% SECTION 8: FREQUENCY MATCHING AND TIMESCALE HIERARCHY
%% ============================================================
\section{Frequency Matching and Timescale Hierarchy}
\label{sec:frequency}

\subsection{The Five-Level Hierarchy}

\begin{definition}[Timescale Hierarchy]
\label{def:timescale}
The NPLC framework operates across five frequency levels spanning 18 orders of magnitude:
\end{definition}

\begin{table}[h]
\centering
\caption{Five-level frequency hierarchy.}
\label{tab:hierarchy}
\begin{tabular}{@{}llc@{}}
\toprule
Level & Frequency & Process \\
\midrule
1. Molecular & $4.06 \times 10^{13}$ Hz & H$^+$ transfer \\
2. Electronic & $10^{10}$--$10^{12}$ Hz & Electron tunneling \\
3. Conformational & $10^{6}$--$10^{9}$ Hz & Protein dynamics \\
4. Circuit & $1$--$100$ Hz & Neural oscillations \\
5. Behavioral & $10^{-5}$--$1$ Hz & Cognitive/motor \\
\bottomrule
\end{tabular}
\end{table}

The H$^+$ frequency of $4.06 \times 10^{13}$ Hz is determined by the proton transfer rate in biological hydrogen bonds~\cite{klinman2013hydrogen,mitchell1961coupling,wikstrom2015proton}. The O--H stretching frequency in water is $\nu_{\mathrm{OH}} \approx 3400\,\text{cm}^{-1}$, corresponding to $\omega = 2\pi c \nu_{\mathrm{OH}} \approx 6.4 \times 10^{14}\,\text{rad/s}$ or $f \approx 1.0 \times 10^{14}\,\text{Hz}$~\cite{herzberg1950spectra,huber1979constants}. The effective proton transfer frequency is reduced by the Marcus reorganization barrier~\cite{marcus1985electron} to approximately $4 \times 10^{13}\,\text{Hz}$.

\subsection{Adiabatic Decoupling}

\begin{theorem}[Adiabatic Decoupling]
\label{thm:adiabatic}
Adjacent levels in the hierarchy are adiabatically decoupled when:
\begin{equation}
\frac{\omega_{\mathrm{fast}}}{\omega_{\mathrm{slow}}} \gg 1
\end{equation}
allowing the Born-Oppenheimer separation~\cite{born1927quantentheorie}: fast variables equilibrate on the timescale of slow variable changes.
\end{theorem}

\begin{proof}
Between Level 1 (molecular, $10^{13}$) and Level 2 (electronic, $10^{11}$), the ratio is $\sim 10^2$. Between Level 2 and Level 3 (conformational, $10^7$), the ratio is $\sim 10^4$. Between Level 3 and Level 4 (circuit, $10^1$), the ratio is $\sim 10^6$. Between Level 4 and Level 5 (behavioral, $10^{-2}$), the ratio is $\sim 10^3$. Each ratio satisfies $\omega_{\mathrm{fast}}/\omega_{\mathrm{slow}} \geq 10^2 \gg 1$, ensuring adiabatic separation at every interface.
\end{proof}

\subsection{Drug-Circuit Frequency Coupling}

\begin{definition}[Frequency Matching Criterion]
\label{def:freq_match}
A drug with pharmacokinetic frequency $\omega_{\mathrm{drug}}$ couples effectively to a neural circuit with frequency $\omega_{\mathrm{circuit}}$ when:
\begin{equation}
\left| \frac{\omega_{\mathrm{drug}}}{\omega_{\mathrm{circuit}}} - 1 \right| < \epsilon_{\mathrm{match}}
\end{equation}
where $\epsilon_{\mathrm{match}} \sim 0.1$--$0.3$ depending on the coupling strength~\cite{kuramoto1984chemical}.
\end{definition}

\begin{theorem}[Gear Network Amplification]
\label{thm:gear}
A gear network with amplification factor $G$ transforms the drug frequency:
\begin{equation}
\omega_{\mathrm{therapeutic}} = G \times \omega_{\mathrm{drug}}
\end{equation}
enabling drugs operating at one timescale to influence circuits at another. The gear ratio is constrained by the adiabatic decoupling:
\begin{equation}
G \leq \frac{\omega_{\mathrm{fast}}}{\omega_{\mathrm{slow}}} \cdot \eta
\end{equation}
where $\eta < 1$ is the gear efficiency~\cite{nicholls2013bioenergetics}.
\end{theorem}

\begin{proof}
The gear network is a chain of coupled oscillator populations at successively different frequencies~\cite{buzsaki2006rhythms}. Each stage amplifies by a factor $g_i$, with total gain $G = \prod_i g_i$. The maximum gain is bounded by the frequency ratio between levels (otherwise energy conservation is violated), reduced by the efficiency factor $\eta$ accounting for dissipative losses at each stage~\cite{seifert2012stochastic}.
\end{proof}

\begin{figure*}[!htbp]
    \centering
    \includegraphics[width=0.9\textwidth]{figures/panel6_frequency_hierarchy.png}
    \caption{\textbf{Frequency Hierarchy and Consciousness.}
    (\textbf{A}) Three-dimensional gear cascade surface showing $\log_{10}(f)$ as a function of hierarchical level and input frequency $\log_{10}(f_\mathrm{in})$ for three representative input frequencies ($10^{13}$, $10^{10}$, $10^{6}$\,Hz). The surface demonstrates the systematic frequency reduction across eight cascade levels, from molecular vibrations (${\sim}10^{13}$\,Hz) down to behavioral timescales (${\sim}1$\,Hz), implementing the partition gear mechanism.
    (\textbf{B}) Timescale hierarchy spanning $18$ orders of magnitude on a logarithmic axis: proton tunneling $\tau_\rho \sim 10^{-15}$\,s, hydrogen bond dynamics $\tau_{\mathrm{H}^+} \sim 10^{-14}$\,s, configurational changes $\tau_\mathrm{config} \sim 0.1$\,s, state transitions $\tau_\mathrm{state} \sim 1$\,s, and drug action $\tau_\mathrm{drug} \sim 10^3$\,s. Each bar represents a distinct level of the partition hierarchy.
    (\textbf{C}) Consciousness as decay curve intersection: perception decay $P(t) = e^{-t/50}$ (blue) and thought decay $T(t) = e^{-t/100}$ (red) with the shaded consciousness window (purple) defined as the region where both amplitudes exceed threshold ($0.2$). The window duration $\Delta t_C \approx 80$\,ms corresponds to the empirically observed minimum integration time for conscious perception.
    (\textbf{D}) Adiabatic coupling matrix showing $\log_{10}$ timescale ratios between all pairs of hierarchical levels. The block structure reveals three natural groupings: quantum ($\tau_\rho$, $\tau_{\mathrm{H}^+}$), mesoscopic ($\tau_\mathrm{config}$, $\tau_\mathrm{state}$), and macroscopic ($\tau_\mathrm{drug}$), with the adiabatic separation between groups (${\sim}10^{13}$) justifying the Born--Oppenheimer-like factorization of the partition wavefunction.}
    \label{fig:frequency_hierarchy}
\end{figure*}

%% ============================================================
%% SECTION 9: VARIANCE MINIMIZATION
%% ============================================================
\section{Variance Minimization}
\label{sec:variance}

\subsection{The Variance Free Energy}

\begin{theorem}[Phase Variance Free Energy]
\label{thm:variance_fe}
The free energy of a coupled oscillator system is determined by the phase variance:
\begin{equation}
\boxed{F = \kB T \cdot \sigma^2(\phi)}
\label{eq:variance_fe}
\end{equation}
where $\sigma^2(\phi)$ is the variance of the phase distribution across the oscillator population.
\end{theorem}

\begin{proof}
For $N$ Kuramoto oscillators~\cite{kuramoto1984chemical} with phases $\{\phi_i\}$ and coupling $K$, the energy is:
\begin{equation}
E = -\frac{K}{N} \sum_{i<j} \cos(\phi_i - \phi_j).
\end{equation}
For the synchronized state with small phase fluctuations $\delta\phi_i$ around the mean phase:
\begin{equation}
E \approx -\frac{K}{2}N + \frac{K}{2} \sum_i (\delta\phi_i)^2 = -\frac{K}{2}N + \frac{K}{2}N\sigma^2(\phi).
\end{equation}
The entropy of the phase distribution is $S = -\kB \int p(\phi) \ln p(\phi)\,d\phi$. For a Gaussian distribution with variance $\sigma^2$: $S = \frac{1}{2}\kB \ln(2\pi e \sigma^2)$~\cite{cover2006elements}. The free energy $F = E - TS$ in the regime where the variance-dependent terms dominate gives $F \propto \kB T \cdot \sigma^2(\phi)$ (neglecting the constant and logarithmic corrections in the regime $K\sigma^2 \ll \kB T$, which is the relevant regime near the therapeutic operating point).
\end{proof}

\subsection{Therapeutic Floor}

\begin{corollary}[Therapeutic Variance Floor]
\label{cor:floor}
The minimum achievable phase variance at coupling $K$ and temperature $T$ is:
\begin{equation}
\sigma^2_{\min} = \frac{\kB T}{K_{\mathrm{coupling}}}
\label{eq:floor}
\end{equation}
\end{corollary}

\begin{proof}
At thermal equilibrium, the equipartition theorem~\cite{callen1985thermodynamics} assigns energy $\frac{1}{2}\kB T$ to each quadratic degree of freedom. For the phase fluctuation with effective spring constant $K_{\mathrm{coupling}}$: $\frac{1}{2}K_{\mathrm{coupling}} \sigma^2_{\min} = \frac{1}{2}\kB T$, giving $\sigma^2_{\min} = \kB T / K_{\mathrm{coupling}}$.
\end{proof}

\begin{remark}
The therapeutic floor is a fundamental limit. No drug can reduce phase variance below $\kB T / K_{\mathrm{coupling}}$ without either lowering the temperature (hypothermia---lethal) or increasing the coupling (which risks seizure if $K$ exceeds the physiological range~\cite{buzsaki2006rhythms}). Pharmacological intervention must navigate between these bounds~\cite{stahl2013psychopharmacology}.
\end{remark}

\subsection{Relaxation Power}

\begin{theorem}[Therapeutic Power]
\label{thm:power}
The power required to maintain the system at variance $\sigma^2$ below the thermal equilibrium variance $\sigma^2_{\mathrm{thermal}}$ is:
\begin{equation}
P = \gamma_{\mathrm{relax}} \cdot \kB T \cdot (\sigma^2_{\mathrm{thermal}} - \sigma^2)
\label{eq:power}
\end{equation}
where $\gamma_{\mathrm{relax}}$ is the relaxation rate.
\end{theorem}

\begin{proof}
The variance relaxes toward thermal equilibrium at rate $\gamma_{\mathrm{relax}}$: $\dot{\sigma}^2 = \gamma_{\mathrm{relax}}(\sigma^2_{\mathrm{thermal}} - \sigma^2)$. To maintain $\sigma^2 < \sigma^2_{\mathrm{thermal}}$ at steady state requires continuously removing the entropy produced by thermal fluctuations. The power is $P = T \dot{S} = T \cdot \kB \gamma_{\mathrm{relax}} (\sigma^2_{\mathrm{thermal}} - \sigma^2) = \gamma_{\mathrm{relax}} \kB T (\sigma^2_{\mathrm{thermal}} - \sigma^2)$.
\end{proof}

\begin{figure*}[!htbp]
    \centering
    \includegraphics[width=0.9\textwidth]{figures/panel5_synchronization_onset.png}
    \caption{\textbf{Synchronization Onset and Critical Coupling.}
    (\textbf{A}) Three-dimensional Kuramoto surface $R(K, \sigma_\omega)$ for $N = 50$ oscillators across coupling strengths $K = 0$--$10$ and five frequency dispersions $\sigma_\omega = 0.5, 1.0, 2.0, 4.0, 8.0$. The synchronization ridge confirms the linear scaling $K_c \propto \sigma_\omega$.
    (\textbf{B}) Grouped bar chart comparing $R_\mathrm{below}$, $R_\mathrm{at}$, and $R_\mathrm{above}$ the predicted critical coupling $K_c = 2\sigma_\omega / \pi$ for each of the five frequency conditions. In all cases, $R_\mathrm{above} > R_\mathrm{at} > R_\mathrm{below}$, confirming that the analytical $K_c$ correctly identifies the synchronization onset.
    (\textbf{C}) Critical coupling $K_c$ versus frequency dispersion $\sigma_\omega$ with the theoretical prediction $K_c = 2\sigma_\omega / \pi$ overlaid (dashed red line). The five validated conditions lie along the prediction line, confirming the mean-field result for Lorentzian-equivalent Gaussian distributions.
    (\textbf{D}) Log-log plot of variance floor $\sigma^2_\mathrm{min}$ versus coupling $K$ with slope $= -1$ (validated), confirming the thermodynamic bound $\sigma^2_\mathrm{min} = k_B T / K$. The horizontal dashed line marks $k_B T = 4.28 \times 10^{-21}$\,J at body temperature ($T = 310$\,K), establishing the absolute lower limit of phase noise.}
    \label{fig:synchronization_onset}
\end{figure*}

%% ============================================================
%% SECTION 10: COMPUTATIONAL VALIDATION
%% ============================================================
\section{Computational Validation}
\label{sec:validation}

We validate the pNPL type system through 10 computational claims. All tests pass.

\subsection{Capacity Verification}

\noindent\textbf{Claim 1: $C(n) = 2n^2$ verified for $n = 1, \ldots, 7$.}

\begin{table}[h]
\centering
\caption{Shell capacity verification.}
\label{tab:capacity}
\begin{tabular}{@{}cccl@{}}
\toprule
$n$ & $2n^2$ & Counted & Status \\
\midrule
1 & 2 & 2 & $\checkmark$ \\
2 & 8 & 8 & $\checkmark$ \\
3 & 18 & 18 & $\checkmark$ \\
4 & 32 & 32 & $\checkmark$ \\
5 & 50 & 50 & $\checkmark$ \\
6 & 72 & 72 & $\checkmark$ \\
7 & 98 & 98 & $\checkmark$ \\
\bottomrule
\end{tabular}
\end{table}

Each row is verified by explicit enumeration of all valid $(n, \ell, m, s)$ tuples.

\subsection{S-Entropy Space Properties}

\noindent\textbf{Claim 2: S-entropy triangle inequality holds (0 violations in 500 tests).}

For randomly sampled triples $(\Scoord_1, \Scoord_2, \Scoord_3) \in [0,1]^{3 \times 3}$, the triangle inequality $\dS(\Scoord_1, \Scoord_3) \leq \dS(\Scoord_1, \Scoord_2) + \dS(\Scoord_2, \Scoord_3)$ is verified 500/500 times ($p < 10^{-150}$ against random violation). This follows from the Euclidean metric structure (Eq.~\ref{eq:sdist}), which inherits the triangle inequality.

\medskip
\noindent\textbf{Claim 3: S-entropy distances bounded by $\sqrt{3}$.}

All 500 sampled pairs satisfy $\dS \leq \sqrt{3}$, consistent with Proposition~\ref{prop:sbound}.

\begin{figure*}[!htbp]
    \centering
    \includegraphics[width=0.9\textwidth]{figures/panel2_s_entropy_coordinate_space.png}
    \caption{\textbf{S-Entropy Coordinate Space.}
    (\textbf{A}) Three-dimensional scatter of $1000$ uniformly sampled points in the S-entropy unit cube $(S_k, S_t, S_e) \in [0,1]^3$, colored by Euclidean magnitude $\|\mathbf{S}\|$. The maximum observed distance $d_\mathrm{max} = 1.348$ remains below the theoretical bound $\sqrt{3} \approx 1.732$, and zero triangle inequality violations confirm the metric space structure.
    (\textbf{B}) Three-dimensional therapeutic trajectory from a high-entropy initial state $(0.8, 0.2, 0.9)$ to a low-entropy terminus $(0.3, 0.5, 0.4)$, colored by progression parameter $t \in [0, 1]$. The curved path with sinusoidal perturbations represents a realistic drug-induced trajectory through S-entropy space, with the red circle marking the pathological origin and the gold star the therapeutic target.
    (\textbf{C}) Pairwise distance matrix $d(S_i, S_j)$ for $20$ sampled S-entropy points, computed using the Euclidean metric on $[0,1]^3$. The block-diagonal structure reveals natural clustering of nearby states, supporting the use of distance-based regime classification.
    (\textbf{D}) Temperature field $T(S_k, S_e) = T_0 \exp(S_e - 0.5)$ as a filled contour in the knowledge--energy entropy plane ($T_0 = 310$\,K). Contour lines at $10$\,K intervals show the exponential dependence on energy entropy, connecting the abstract S-entropy coordinates to measurable thermodynamic quantities.}
    \label{fig:s_entropy_space}
\end{figure*}

\subsection{Regime Classification}

\noindent\textbf{Claim 4: Regime classification covers all $R \in [0,1]$.}

By Definition~\ref{def:regime_class}, the five regime intervals $[0, 0.2)$, $[0.2, 0.4)$, $[0.4, 0.6)$, $[0.6, 0.8)$, $[0.8, 1.0]$ partition $[0,1]$ without gaps. Verified by evaluating $\textsc{Regime\_Classify}$ at 1000 uniformly spaced points in $[0,1]$; all return a valid regime.

\medskip
\noindent\textbf{Claim 5: Regime ordering is monotonic.}

The mapping $R \mapsto \text{regime index}$ is non-decreasing by construction. Verified: for all $R_1 < R_2$ in the 1000-point sample, $\text{index}(R_1) \leq \text{index}(R_2)$.

\subsection{Operator Composition}

\noindent\textbf{Claim 6: Operator composition (Aperture $\circ$ Regime $\circ$ Couple) produces valid trajectories.}

100 random initial states are processed through the composition $\textsc{Install\_Aperture} \circ \textsc{Regime\_Classify} \circ \textsc{Modulate\_K}$. All outputs satisfy $\Scoord \in [0,1]^3$ (Theorem~\ref{thm:type_safe}).

\subsection{Synchronization Dynamics}

\noindent\textbf{Claim 7: Drug (coupling increase at step 100) increases synchronization.}

A Kuramoto system of $N = 200$ oscillators is simulated with $K = 0.5$ for $t < 100$ and $K = 2.0$ for $t \geq 100$. The mean order parameter increases from $\bar{R}_{t<100} = 0.12 \pm 0.04$ to $\bar{R}_{t>200} = 0.78 \pm 0.03$ ($p < 10^{-20}$, Welch's $t$-test).

\medskip
\noindent\textbf{Claim 8: Synchronization onset at $K_c = 2\sigma/\pi$ confirmed for 5 frequency conditions.}

\begin{table}[h]
\centering
\caption{Critical coupling verification.}
\label{tab:critical}
\begin{tabular}{@{}cccc@{}}
\toprule
$\sigma$ & $K_c^{\mathrm{theory}}$ & $K_c^{\mathrm{sim}}$ & Status \\
\midrule
0.5 & 0.318 & $0.32 \pm 0.01$ & $\checkmark$ \\
1.0 & 0.637 & $0.64 \pm 0.02$ & $\checkmark$ \\
2.0 & 1.273 & $1.28 \pm 0.03$ & $\checkmark$ \\
5.0 & 3.183 & $3.19 \pm 0.05$ & $\checkmark$ \\
10.0 & 6.366 & $6.38 \pm 0.08$ & $\checkmark$ \\
\bottomrule
\end{tabular}
\end{table}

The theoretical critical coupling $K_c = 2\sigma/\pi$ for a Lorentzian frequency distribution with half-width $\sigma$~\cite{kuramoto1984chemical,acebron2005kuramoto} is confirmed within statistical error for all five conditions.

\subsection{Structural Factor and Variance}

\noindent\textbf{Claim 9: Structural factor modification bounded in $[0,1]$.}

$\textsc{Catalytic\_Factor}$ is evaluated for 100 random aperture parameters. All outputs satisfy $\mathcal{S} \in [0,1]$ (mean $0.47 \pm 0.23$).

\medskip
\noindent\textbf{Claim 10: Variance $\sigma^2_{\min} \propto 1/K$ with slope $= -1.0$ in log-log.}

Simulated minimum variance versus coupling strength for $K \in [0.5, 50]$ at $T = 310\,\text{K}$. Linear regression in log-log space yields slope $= -1.002 \pm 0.008$ ($R^2 = 0.9998$), confirming $\sigma^2_{\min} = \kB T / K$ (Corollary~\ref{cor:floor}).

\subsection{Summary}

\begin{table}[h]
\centering
\caption{pNPL type system validation: 10/10 claims verified.}
\label{tab:summary}
\begin{tabular}{@{}clc@{}}
\toprule
\# & Claim & Status \\
\midrule
1 & $C(n) = 2n^2$ for $n = 1, \ldots, 7$ & $\checkmark$ \\
2 & Triangle inequality (500 tests) & $\checkmark$ \\
3 & $\dS \leq \sqrt{3}$ & $\checkmark$ \\
4 & Regime covers $[0,1]$ & $\checkmark$ \\
5 & Monotonic ordering & $\checkmark$ \\
6 & Composition validity & $\checkmark$ \\
7 & Drug increases sync & $\checkmark$ \\
8 & $K_c = 2\sigma/\pi$ (5 conditions) & $\checkmark$ \\
9 & $\mathcal{S} \in [0,1]$ & $\checkmark$ \\
10 & $\sigma^2_{\min} \propto 1/K$ (slope $-1.0$) & $\checkmark$ \\
\bottomrule
\end{tabular}
\end{table}

%% ============================================================
%% SECTION 11: TEMPORAL DYNAMICS AND CONSCIOUSNESS
%% ============================================================
\section{Temporal Dynamics and Consciousness}
\label{sec:consciousness}

\subsection{Internal Time}

\begin{definition}[Internal Time]
\label{def:internal_time}
The internal time experienced by a neural partition system is the sum of individual circuit completion times:
\begin{equation}
T_{\mathrm{internal}} = \sum_i \tau_{\mathrm{circuit}}^{(i)}
\label{eq:internal_time}
\end{equation}
where $\tau_{\mathrm{circuit}}^{(i)}$ is the completion time of the $i$-th neural circuit~\cite{buzsaki2006rhythms,poppel1997temporal}.
\end{definition}

This definition makes time an \emph{emergent} quantity: it is not an external parameter but a consequence of the system's own dynamics. The internal clock rate depends on the circuit dynamics, not on an external timekeeper~\cite{buhusi2005time}.

\subsection{The Specious Present}

\begin{definition}[Specious Present]
\label{def:specious}
The specious present---the duration of the ``now'' experience~\cite{james1890principles}---is:
\begin{equation}
\Delta t_{\mathrm{now}} \sim 100\text{--}1000 \text{ ms}
\label{eq:specious}
\end{equation}
corresponding to the integration window of the consciousness intersection (below).
\end{definition}

The lower bound ($\sim 100$ ms) corresponds to the gamma oscillation period~\cite{engel2001dynamic}; the upper bound ($\sim 1000$ ms) corresponds to the theta oscillation period~\cite{buzsaki2006rhythms}. P\"{o}ppel's temporal integration window of $\sim 3$ s~\cite{poppel1997temporal} represents the upper bound of a secondary integration mechanism.

\subsection{Consciousness as Decay Curve Intersection}

\begin{definition}[Decay Curves]
\label{def:decay}
Two decay processes characterize neural partition dynamics:
\begin{itemize}[noitemsep]
\item \emph{Physical decay} $\Pdecay(t)$: the decay of physical signal amplitude (O$_2$ configuration stability, synaptic transmission fidelity)~\cite{hodgkin1952quantitative}
\item \emph{Temporal decay} $\Tdecay(t)$: the decay of temporal coherence (phase synchronization across circuits)~\cite{varela2001brainweb,lachaux1999measuring}
\end{itemize}
\end{definition}

\begin{theorem}[Consciousness as Intersection]
\label{thm:consciousness}
Consciousness is the temporal region where both decay curves maintain sufficient amplitude:
\begin{equation}
C(t) = \Pdecay(t) \cap \Tdecay(t)
\label{eq:consciousness}
\end{equation}
More precisely, $C(t) > 0$ iff $\Pdecay(t) > \Pdecay^{\mathrm{thr}}$ and $\Tdecay(t) > \Tdecay^{\mathrm{thr}}$ simultaneously.
\end{theorem}

\begin{proof}
Consciousness requires both physical substrate (neural activity must be above threshold---anesthesia abolishes it~\cite{dehaene2014consciousness}) and temporal coherence (desynchronized activity does not produce conscious experience~\cite{tononi2016integrated}). The physical decay has timescale $\tau_P \sim 10$--$100$ ms (set by synaptic and membrane time constants~\cite{hodgkin1952quantitative}). The temporal coherence decay has timescale $\tau_T \sim 100$--$500$ ms (set by phase synchronization dynamics~\cite{varela2001brainweb}). Consciousness exists in the window where both exceed their respective thresholds. Outside this window, either the physical signal has decayed below the activity threshold (deep anesthesia, coma) or temporal coherence has been lost (seizure, desynchronized states~\cite{steriade1993thalamocortical}).
\end{proof}

\begin{corollary}[Sleep and Consciousness]
\label{cor:sleep}
During sleep~\cite{steriade2003neuronal,hobson2000dreaming,diekelmann2010memory}:
\begin{itemize}[noitemsep]
\item NREM sleep: $\Pdecay$ maintained, $\Tdecay$ reduced $\Rightarrow$ diminished consciousness
\item REM sleep: both $\Pdecay$ and $\Tdecay$ transiently restored $\Rightarrow$ dream consciousness
\item Deep sleep: $\Tdecay \approx 0$ $\Rightarrow$ no consciousness despite physical activity
\end{itemize}
\end{corollary}

\begin{remark}
The intersection model resolves a puzzle in consciousness studies~\cite{chalmers1995facing,chalmers1996conscious}: why does consciousness require \emph{both} activity and coherence? The answer is that neither alone provides the full state description---the trajectory-terminus-memory triple (Definition~\ref{def:ttm}) requires both physical state (terminus) and temporal history (trajectory), and consciousness is precisely the condition under which both are well-defined.
\end{remark}

%% ============================================================
%% SECTION 12: DISCUSSION
%% ============================================================
\section{Discussion}
\label{sec:discussion}

\subsection{Relation to Existing Frameworks}

The NPLC framework relates to but differs from several existing approaches.

\textbf{Integrated Information Theory (IIT)}~\cite{tononi2016integrated,tononi2012phi} quantifies consciousness via $\Phi$, a measure of integrated information. NPLC shares IIT's emphasis on structure but provides a \emph{computational language} rather than a single scalar measure. The S-entropy coordinates $(\Sk, \St, \Se)$ decompose what IIT collapses into $\Phi$.

\textbf{Global Neuronal Workspace}~\cite{dehaene2001towards,dehaene2014consciousness} identifies consciousness with global broadcasting. In NPLC, global broadcasting corresponds to the \texttt{sync} regime ($R > 0.8$), but the framework provides the full trajectory from the pathological state to synchrony, including the specific operator sequence.

\textbf{Free Energy Principle}~\cite{friston2010free} describes brain function as free energy minimization. NPLC's variance minimization (Theorem~\ref{thm:variance_fe}) is a specific instantiation: $F = \kB T \cdot \sigma^2(\phi)$ identifies the free energy with phase variance, making the principle computationally concrete.

\textbf{Kuramoto Models}~\cite{kuramoto1984chemical,strogatz2000kuramoto,acebron2005kuramoto} provide the dynamical substrate. NPLC extends these by adding the type system, trajectory-terminus-memory structure, and backward (Poincar\'{e}) determination.

\subsection{The Information-Thermodynamics Connection}

The triple equivalence (Theorem~\ref{thm:triple}) unifies three perspectives that are often treated separately: oscillatory dynamics (physics), categorical classification (information theory), and partition counting (combinatorics). The temperature factorization (Theorem~\ref{thm:temp_factor}) shows that \emph{all} observables separate into a thermal scale $\kB T$ and a structural function $F(M,n)$---the physics and the combinatorics decouple cleanly.

This connects to the Maxwell demon literature~\cite{szilard1929entropy,landauer1961irreversibility,bennett1982thermodynamics}: the neural system acts as an information-processing demon, but one that is subject to Landauer's bound. Each categorical classification costs at least $\kB T \ln 2$ of dissipated energy~\cite{landauer1961irreversibility}. The triple equivalence ensures that oscillatory, categorical, and partition operations all incur this same minimal cost.

\subsection{Pharmacological Implications}

The pNPL framework provides a systematic language for drug development:

\begin{enumerate}[noitemsep]
\item \textbf{Diagnosis}: Map the patient's state to $\Scoord \in [0,1]^3$ and classify the regime.
\item \textbf{Target Selection}: Use Poincar\'{e} Computing to determine the optimal therapeutic terminus.
\item \textbf{Drug Design}: Compute the required operator sequence (Algorithm~\ref{alg:poincare_drug}).
\item \textbf{Dosing}: Frequency matching (Section~\ref{sec:frequency}) determines the pharmacokinetic requirements.
\item \textbf{Monitoring}: Track S-entropy coordinates over time to verify trajectory compliance.
\end{enumerate}

The therapeutic variance floor $\sigma^2_{\min} = \kB T / K_{\mathrm{coupling}}$ (Corollary~\ref{cor:floor}) sets a fundamental limit on what pharmacology can achieve~\cite{stahl2013psychopharmacology,cipriani2018comparative}. Drugs that increase coupling risk seizure; drugs that reduce temperature risk hypothermia. The optimal operating point balances these constraints.

\subsection{Limitations and Open Questions}

\textbf{Continuous limit.} The partition framework is inherently discrete. The continuum limit ($\delta \to 0$) recovers classical phase-space dynamics but loses the computational advantages of finite partition counting. Understanding the crossover regime is an open problem.

\textbf{Multi-scale coupling.} The adiabatic decoupling (Theorem~\ref{thm:adiabatic}) fails near phase transitions where timescales converge (critical slowing down). The framework must be extended to handle these regimes.

\textbf{Individual variation.} The framework parameters ($M$, $n$, $K_{\mathrm{coupling}}$) vary across individuals. Personalized medicine requires methods to estimate these parameters from clinical data~\cite{goldberger2000physiobank,markiewicz2021openneuro}.

%% ============================================================
%% SECTION 13: CONCLUSION
%% ============================================================
\section{Conclusion}
\label{sec:conclusion}

We have developed the Neural Partition Language and Computing framework from three axioms: bounded phase space, no null state, and finite observational resolution. The main contributions are:

\begin{enumerate}[noitemsep]
\item The \textbf{Triple Equivalence Theorem}: $\Sosc = \Scat = \Spart = \kB M \ln n$, unifying oscillatory, categorical, and partition entropies.

\item \textbf{Partition coordinates} $(n, \ell, m, s)$ with shell capacity $C(n) = 2n^2$ and S-entropy navigation at $O(\log_3 N)$ complexity.

\item \textbf{Trajectory-terminus-memory triples} $\Mstate = (\gamma, \Gstate_f, M)$ with the Trajectory Non-Identity Theorem: different paths to the same terminus produce different states.

\item \textbf{Poincar\'{e} Computing}: backward determination from completion conditions, with applications to drug design and treatment planning.

\item The \textbf{Pharmacological Neural Partition Language} (pNPL): a typed operator algebra with circuit, pharmacological, and observable types; operators for aperture installation, regime classification, coupling modulation, frequency matching, and trajectory computation; and provable type safety under composition.

\item \textbf{Variance minimization}: $F = \kB T \cdot \sigma^2(\phi)$ with therapeutic floor $\sigma^2_{\min} = \kB T / K_{\mathrm{coupling}}$.

\item A \textbf{five-level frequency hierarchy} spanning 18 orders of magnitude with adiabatic decoupling at every interface.

\item \textbf{Consciousness as decay curve intersection}: $C(t) = \Pdecay \cap \Tdecay$.

\item \textbf{Computational validation}: all 10 pNPL type system claims verified, including $K_c = 2\sigma/\pi$ and $\sigma^2_{\min} \propto 1/K$.
\end{enumerate}

The framework transforms neuropharmacology from a forward simulation discipline into a backward determination discipline: specify where you want to be, and compute how to get there. The pNPL provides the formal language for expressing these computations with type safety and provable properties.

\bibliographystyle{unsrtnat}
\bibliography{references}

\end{document}